\section*{Заключение}
\addcontentsline{toc}{section}{Заключение}

В ходе выполнения курсовой работы были решены следующие задачи:

\begin{enumerate}
    \item Изучены теоретические основы гидродинамической смазки и влияние дискретного микрорельефа на трибологические характеристики подшипников скольжения.

    \item Освоена методика численного решения уравнения Рейнольдса методом конечных разностей с итерационной процедурой Гаусса--Зейделя.

    \item Выполнен расчёт поля давления и интегральных характеристик радиального подшипника скольжения для гладкой поверхности и поверхности с \varDimpleTypeGenitive{} углублениями (тип~\varDimpleType, вариант~\varNumber).

    \item Проведён сравнительный анализ результатов в диапазоне эксцентриситетов $\varepsilon = 0.05 \div 0.8$.
\end{enumerate}

Полученные результаты подтверждают эффективность применения дискретного микрорельефа для улучшения характеристик подшипников скольжения. Увеличение несущей способности обусловлено микрогидродинамическим эффектом, а снижение коэффициента трения~--- увеличением средней толщины смазочного слоя в зоне углублений.

\newpage
